\documentclass{article}

\usepackage[preprint]{neurips_2024}

% =============================================================================
% Packages
% =============================================================================

\usepackage[T1]{fontenc}
\usepackage[utf8]{inputenc}
\usepackage{times}
\usepackage{microtype}
\usepackage{inconsolata}

\usepackage{amsmath}
\usepackage{amssymb}

\usepackage{graphicx}
\usepackage{booktabs}
\usepackage{multirow}

\usepackage{hyperref}
\usepackage{xurl}
\usepackage{placeins}
\usepackage{float}
\usepackage{enumitem}
\usepackage{algorithm}
\usepackage{algorithmic}

% =============================================================================
% Title
% =============================================================================

\title{Typed Composition Routing:\\Tool Selection via Entity Type Decomposition}

\author{
  Jose Angel Morena Simon
}

\begin{document}

\maketitle

% =============================================================================
% Abstract
% =============================================================================

\begin{abstract}
As LLM agents integrate with increasingly large tool ecosystems, tool routing---the task of determining which tools to invoke for a given query---becomes a major bottleneck. Existing approaches require the model to simultaneously identify relevant tools, infer valid multi-step compositions, and avoid hallucinated tool calls. We propose Typed Composition Routing (TCR), which decomposes this problem into two stages: the model predicts source and target entity types, while deterministic graph search over a typed composition graph reconstructs valid tool chains. Across four LLMs, five domains (54--170 tools), and 140 queries, TCR improves F1 in all 20 model--domain combinations (average F1 $+0.32$). Because the graph contains only registered tools, nonexistent tool invocations are eliminated by construction. On a production MCP server with 1,060 tools, TCR reduces prompt tokens by 99.7\%. Finally, a granularity study across two domains shows that type ontology design is a key performance lever whose optimal granularity depends on the capability of the entity-type predictor.
\end{abstract}

% =============================================================================
% Introduction
% =============================================================================

\section{Introduction}

LLM agents increasingly rely on external tools, but existing approaches typically expose the entire tool catalog to the model. As catalogs grow from tens to thousands of APIs, this increases prompt tokens, inference cost, and routing difficulty.

Existing approaches leave both semantic and compositional reasoning entirely to the language model, without exploiting the typed structure already present in many APIs.

A simple observation motivates a different approach. Users naturally express intent in terms of domain entities rather than tools. A query like ``get the logs for pods in production'' refers to entity types---Namespace, PodLogs---without naming any tool. Given a typed API specification, valid tool chains can be recovered deterministically from the type graph. Entity types are explicitly signaled by user intent, whereas the tools that connect them are latent variables that must be inferred. By targeting entity types, the model classifies what the user already said rather than inferring what the user never mentioned.

\textbf{Typed Composition Routing (TCR)} addresses this by decoupling semantic reasoning from compositional reasoning. The model handles only (1)~\textbf{semantic reasoning}---predicting source and target entity types---while (2)~\textbf{compositional reasoning} is delegated to deterministic graph search over a typed composition graph, which reconstructs structurally valid tool chains. The composition graph is derived from typed API specifications, automatically when schemas expose input/output types (e.g., OpenAPI), or through lightweight manual annotations otherwise.

We make three contributions:
\begin{enumerate}
  \item A method that decomposes tool selection into entity-type prediction and deterministic graph search, providing compositional validity guarantees and eliminating nonexistent tool invocations by construction.
  \item Experiments across four LLMs and five domains demonstrate consistent improvements over direct selection (average F1 $+0.32$). On a production MCP server with 1,060 tools, TCR reduces prompt tokens by 99.7\% while eliminating hallucinated tool invocations.
  \item A study of ontology granularity across two domains showing that type granularity is an important design parameter whose optimal value depends on predictor capacity.
\end{enumerate}

% =============================================================================
% Method
% =============================================================================

\section{Method}

\subsection{Typed Composition Graph}

We model each tool as a typed transformation: a tool $t$ takes an entity of type $e_i$ and produces an entity of type $e_j$. The full tool registry then induces a directed graph $G = (E, T)$, where the $|E|$ entity types form the nodes and the $|T|$ tools induce directed edges between compatible types. Any valid multi-step workflow corresponds to a path through this graph.

\subsection{Routing Decomposition}

Given a natural language query $q$, TCR proceeds in two stages. First, a predictor identifies the entity types the user is referring to:
\begin{equation}
  f(q) \rightarrow (e_{\text{src}}, e_{\text{tgt}})
\end{equation}
Then, deterministic graph search enumerates all type-compatible tool chains connecting those types:
\begin{equation}
  \text{Compose}(G, e_{\text{src}}, e_{\text{tgt}}) \rightarrow \{[t_1, \ldots, t_l]\}
\end{equation}

This decomposition rests on an asymmetry: entity types are explicit in user queries (``pods,'' ``credentials,'' ``inventories''), whereas tool identities and their compositions are latent. This reduces the prediction space from hundreds or thousands of tools to a much smaller set of semantic types---the model classifies what the user already said rather than inferring what the user never mentioned. Because the graph contains only registered tools as edges, any traversal inherits a structural guarantee: every returned tool belongs to the registry and every adjacent pair is type-compatible by construction, unlike standard retrieval or direct LLM selection approaches, which do not enforce these guarantees. The candidate set can be used directly (when it is small enough) or passed to a second-stage LLM that selects from the narrowed list rather than the full catalog. This allows a downstream model to reason over a dramatically smaller search space without sacrificing valid compositions. Importantly, the graph search stage performs no language understanding or semantic inference; all semantic reasoning is confined to entity-type prediction.

\textbf{Example.} Consider the query \textit{``Get the logs for pods in the production namespace.''}  The predictor identifies Namespace as the source type and PodLogs as the target. Graph search finds the path Namespace $\rightarrow$ Pod $\rightarrow$ PodLogs, yielding the tool chain [\texttt{list\_namespaced\_pods}, \texttt{get\_pod\_logs}]---two tools the user never named. Direct selection would need to identify both tools and their ordering; TCR predicts only the chain endpoints and lets the graph fill in the rest. The predictor never needs to infer the intermediate Pod entity---it emerges naturally from graph traversal. Figure~\ref{fig:method} illustrates the full pipeline.

\begin{figure}[htbp]
  \centering
  \includegraphics[width=\textwidth]{figures/1_method.pdf}
  \caption{The TCR pipeline. A learned predictor maps the query to source and target entity types (Stage~1); deterministic typed graph search produces a filtered tool set of structurally valid candidates (Stage~2); only this filtered set is presented to the LLM for final tool selection (Stage~3).}
  \label{fig:method}
\end{figure}

\subsection{Candidate Search Algorithm}

$\text{Compose}(G, e_{\text{src}}, e_{\text{tgt}})$ returns all tools that lie on type-compatible paths between source and target. The search must satisfy three properties: (1)~every returned tool belongs to at least one valid path from $e_{\text{src}}$ to $e_{\text{tgt}}$; (2)~every returned tool belongs to the registered tool catalog; (3)~self-loop tools on $e_{\text{tgt}}$ (e.g., delete or update operations on the target entity) are included even though they do not advance along a path.
Candidate search performs no semantic reasoning; it operates purely over the typed composition graph.

Exact enumeration of all path-participating edges is unnecessary for our setting. Algorithm~\ref{alg:candidate} describes a linear-time reachability-based approximation that produces exact results on K8s (21 types) and a near-exact superset on AAP (88 types, 3 extra tools on one type pair).

\begin{algorithm}[h]
\caption{CandidateSearch$(G, e_{\text{src}}, e_{\text{tgt}})$}
\label{alg:candidate}
\begin{algorithmic}[1]
\IF{$e_{\text{src}} = e_{\text{tgt}}$}
  \RETURN self-loop tools on $e_{\text{src}}$
\ENDIF
\STATE $G' \leftarrow G$ with self-loop edges removed \COMMENT{they preserve the current entity type}
\STATE $F \leftarrow$ forward-reachable nodes from $e_{\text{src}}$ in $G'$
\STATE $B \leftarrow$ backward-reachable nodes from $e_{\text{tgt}}$ in $G'$
\STATE $V \leftarrow F \cap B$
\STATE $C \leftarrow \{t \mid (u \xrightarrow{t} v) \in G',\; u \in V,\; v \in V\}$
\STATE $C \leftarrow C \;\cup$ self-loop tools on $e_{\text{tgt}}$
\RETURN $C$
\end{algorithmic}
\end{algorithm}

\section{Related Work}

Existing tool-routing methods fall into three broad categories. \textit{Direct selection} approaches train or prompt LLMs to choose tools from catalogs \citep{qin2024toolllm,patil2024gorilla,schick2023toolformer,du2024anytool}. \textit{Graph-constrained planning} introduces structural guidance---for example, ControlLLM~\citep{liu2024controlllm} uses a bipartite resource--tool graph for multimodal routing, while ToolNet~\citep{liu2024toolnet} and GNN4Plan~\citep{wu2024gnn4plan} leverage co-occurrence and graph neural networks to improve tool selection. \textit{Agent frameworks} such as HuggingGPT~\citep{shen2023hugginggpt} and ReAct~\citep{yao2023react} instead rely on open-ended LLM reasoning without explicit structural guarantees. Despite these differences, these methods ultimately formulate tool routing as a direct tool prediction problem.

In contrast, TCR formulates tool routing differently. Rather than predicting tools directly, it predicts semantic endpoint types and delegates tool composition to deterministic typed graph search, thereby decoupling semantic prediction from structural reasoning.

% =============================================================================
% Experiments
% =============================================================================

\section{Experiments}

\subsection{Setup}

We evaluate on five benchmark domains---Kubernetes (135 tools), Ansible (108), GitHub (133), CI/CD (54), and Shopify (170)---plus a production-scale domain: the Ansible Automation Platform (AAP) MCP Server~\citep{anthropic2024mcp}, with 1,060 tools and 50 entity types.

Four LLMs serve as type predictors: Granite 4.1 8B (IBM), Qwen3-14B (Alibaba), GPT-OSS-20B \citep{openai2025gptoss}, and Claude Haiku 4.5 (Anthropic). In the main experiments (\S\ref{sec:main}--\S\ref{sec:production}), all models are evaluated zero-shot, without fine-tuning; the granularity study (\S\ref{sec:granularity}) uses 10 few-shot examples per granularity level to enable comparison across ontology sizes. We compare against two baselines---direct selection (full tool catalog in the prompt) and embedding retrieval (top-10 tools by cosine similarity using nomic-embed-text-v1.5)---as well as an oracle upper bound that bypasses the predictor and uses the ground-truth source and target entity types. All strategies are evaluated as candidate generators: each produces a set of tools for a query, and we report per-query precision (fraction of returned tools in the gold set), recall (fraction of gold tools returned), and their harmonic mean F1, macro-averaged across queries.

The benchmark comprises 140 manually authored queries for the five benchmark domains and 120 queries for AAP, each spanning six categories: \emph{clean} (direct, unambiguous requests), \emph{synonym} (paraphrased intents using different vocabulary), \emph{multi-hop} (queries requiring two or more chained tools), \emph{ambiguous} (underspecified queries with implicit context), \emph{noisy} (conversational or informal phrasing), and \emph{multi-path} (queries where multiple valid tool paths exist). This category mix tests robustness across the variation patterns that arise in real user inputs.

All queries and gold annotations were authored by the paper's authors: queries were written from API documentation without reference to the routing system, expected tool chains were determined by manual inspection of the type graph, and no annotations were revised after observing system outputs. We acknowledge that author-authored benchmarks may unintentionally favor the proposed decomposition; independently constructed benchmarks would strengthen the evaluation.

\subsection{Main Results}\label{sec:main}

Table~\ref{tab:main} summarizes the results. TCR improves F1 in every one of the 20 model--domain combinations, with an average absolute F1 improvement of $+0.32$ (median $+0.31$, SD $0.12$) over direct selection. The effect is statistically significant (Wilcoxon signed-rank test, $T{=}0$, $p = 8.83 \times 10^{-5}$, $N{=}20$) and practically large: even the weakest model under TCR (GPT-OSS-20B at 0.83) outperforms the strongest baseline (Qwen3-14B at 0.58). Because graph search returns only registered tools, TCR cannot produce nonexistent tool invocations---a structural guarantee confirmed empirically across all benchmark queries. The oracle represents an upper bound obtained by providing the true source and target entity types. Since the benchmark tool chains were annotated directly from the typed composition graph, the oracle recovers every gold chain (F1\,=\,1.00). This measures graph completeness rather than predictor performance.

\begin{table}[htbp]
  \centering
  \caption{Average macro-F1 across the five benchmark domains. TCR guarantees zero nonexistent tool invocations by construction.}
  \label{tab:main}
  \begin{tabular}{lcccc}
    \toprule
    \textbf{Model} & \textbf{Baseline} & \textbf{Retrieval} & \textbf{TCR} & \textbf{Oracle} \\
    \midrule
    Granite 4.1 8B  & 0.49 & 0.47 & \textbf{0.85} & 1.00 \\
    Qwen3-14B       & 0.58 & 0.58 & \textbf{0.84} & 1.00 \\
    GPT-OSS-20B     & 0.45 & 0.53 & \textbf{0.83} & 1.00 \\
    Haiku 4.5       & 0.56 & 0.57 & \textbf{0.84} & 1.00 \\
    \bottomrule
  \end{tabular}
\end{table}

\begin{table}[htbp]
  \centering
  \caption{Per-domain F1 for all model--domain combinations. Absolute F1 values are reported as Baseline $\rightarrow$ TCR.}
  \label{tab:perdomain}
  \small
  \begin{tabular}{lccccc}
    \toprule
    \textbf{Model} & \textbf{K8s} & \textbf{Ansible} & \textbf{GitHub} & \textbf{CI/CD} & \textbf{Shopify} \\
    \midrule
    Granite 4.1 8B  & 0.53$\to$0.86 & 0.34$\to$0.81 & 0.36$\to$0.90 & 0.51$\to$0.78 & 0.69$\to$0.90 \\
    Qwen3-14B       & 0.60$\to$0.91 & 0.41$\to$0.71 & 0.56$\to$0.82 & 0.51$\to$0.84 & 0.79$\to$0.92 \\
    GPT-OSS-20B     & 0.55$\to$0.83 & 0.31$\to$0.76 & 0.45$\to$0.89 & 0.54$\to$0.79 & 0.39$\to$0.90 \\
    Haiku 4.5       & 0.58$\to$0.91 & 0.30$\to$0.75 & 0.59$\to$0.88 & 0.55$\to$0.73 & 0.80$\to$0.95 \\
    \bottomrule
  \end{tabular}
\end{table}

The largest gains appear where baseline performance is weakest---on Ansible, per-model F1 improvements range from $+0.30$ (Qwen3) to $+0.47$ (Granite) across the four predictors (Table~\ref{tab:perdomain})---suggesting that typed graph constraints provide the greatest benefit when unconstrained tool selection is more difficult. Separate experiments with Gemini 2.5 Pro on three additional domains (Table~\ref{tab:gemini}) illustrate TCR's behavior with a frontier model. Gemini achieves strong direct-selection F1 (0.69--0.86), so TCR's primary value here is not F1 improvement but structural guarantees: high-recall candidate sets (0.88--0.91) with zero hallucinated tool invocations and substantially reduced prompt tokens. The lower TCR F1 reflects an evaluation artifact: because the benchmark annotates a single gold tool chain, candidate sets containing additional structurally valid tools are penalized as false positives. This precision--recall tradeoff is inherent to candidate-generation evaluation and affects high-recall systems disproportionately. Notably, even Gemini's baseline F1 degrades from 0.86 (17 tools) to 0.69 (1,060 tools), illustrating the scalability challenge that motivates TCR---a challenge that is more acute for the smaller models evaluated in Table~\ref{tab:main}. Overall, these results support the central hypothesis of TCR: decoupling semantic prediction from compositional reasoning improves routing quality, while the structural properties of graph search guarantee the absence of hallucinated tool invocations.

\begin{table}[htbp]
  \centering
  \caption{Gemini 2.5 Pro on three additional domains. TCR candidate sets maintain high recall with zero hallucinated tool invocations. Baseline and TCR columns report F1, matching Table~\ref{tab:main}.}
  \label{tab:gemini}
  \begin{tabular}{lrcccc}
    \toprule
    \textbf{Domain} & \textbf{Tools} & \textbf{Queries} & \textbf{Baseline F1} & \textbf{TCR F1} & \textbf{TCR Recall} \\
    \midrule
    Kubernetes  & 17    & 50  & 0.86 & 0.76 & 0.91 \\
    Stripe      & 536   & 50  & 0.84 & 0.59 & 0.88 \\
    AAP         & 1,060 & 120 & 0.69 & 0.55 & 0.91 \\
    \bottomrule
  \end{tabular}
\end{table}

\subsection{Production Scale: 1,060 Tools}\label{sec:production}

At production scale, direct tool selection incurs substantial prompt overhead. Presenting all 1,060 tool schemas to the LLM requires roughly 25,000 prompt tokens per query. TCR avoids prompting with the full tool catalog: the model receives only entity type names (${\sim}$66 tokens), and graph search generates candidates from the composition graph.

\begin{table}[htbp]
  \centering
  \caption{Production-scale routing on the AAP MCP Server (1,060 tools). \textsuperscript{$\dagger$}Routing overhead measures the additional latency introduced by TCR's entity-type prediction call; direct selection has no separate routing step.}
  \label{tab:production}
  \begin{tabular}{lcccc}
    \toprule
    \textbf{Strategy} & \textbf{Recall} & \textbf{Prompt Tok.} & \textbf{Hallucinations} & \textbf{Routing Overhead}\textsuperscript{$\dagger$} \\
    \midrule
    Direct selection         & 0.58          & 25,433  & 10  & 0\,ms      \\
    TCR + Granite 8B         & 0.42          & 66      & 0   & 6,579\,ms  \\
    TCR + Opus 4.6            & \textbf{1.00} & \textbf{66} & \textbf{0} & 2,100\,ms  \\
    \bottomrule
  \end{tabular}
\end{table}

Table~\ref{tab:production} shows the comparison. TCR replaces a 25,433-token tool catalog with a 66-token entity-type prompt---a 99.7\% reduction. The average candidate set shrinks to 6.1 tools, a 99.4\% reduction from 1,060. Direct selection produced 10 hallucinated tool invocations across 120 queries, whereas TCR produced zero, reflecting TCR's structural guarantee that only registered tools can be returned. Recall depends on predictor capacity: Granite 8B achieves 0.42 recall (primarily limited by its 61\% entity-type prediction accuracy at 50 types), while Opus 4.6 achieved 100\% type-prediction accuracy on this benchmark and 1.00 recall---outperforming direct selection in both recall and hallucinated tool invocations while requiring only 66 prompt tokens. Latency is dominated by the entity-type prediction LLM call rather than graph search, whose runtime is negligible. Figure~\ref{fig:scale} illustrates the simultaneous reduction in candidate set size and prompt tokens enabled by TCR. More broadly, TCR shifts the scaling bottleneck from tool catalogs to semantic type prediction: prompt size grows with the number of entity types rather than the number of registered tools.

\begin{figure}[t!]
  \centering
  \includegraphics[width=\textwidth]{figures/2_scale.pdf}
  \caption{Production-scale routing on the AAP MCP Server. TCR reduces the average candidate set from 1,060 tools to 6.1 (99.4\%) and the prompt from 25,433 tokens to 66 entity-type tokens (99.7\%).}
  \label{fig:scale}
\end{figure}

\FloatBarrier
\subsection{Ontology Granularity Experiment}\label{sec:granularity}

A key design question is how fine-grained the entity ontology should be. The results above treat the type system as fixed. But how sensitive is routing quality to the granularity of the ontology? To find out, we constructed three ontology variants---coarse, medium, and fine---on both a K8s subset (17 tools, 128 queries) and AAP (1,060 tools, 120 queries), varying only the type system while holding tools and queries constant. On K8s, coarse uses 6 entity types (e.g., \texttt{Pod}, \texttt{Resource}), medium splits these into 13 types (distinguishing \texttt{Deployment}, \texttt{Service}, \texttt{Node}, etc.), and fine introduces 21 types with detail variants (\texttt{PodDetail}, \texttt{ServiceDetail}, \texttt{PodCreation}). On AAP, the three levels have 8, 50, and 88 types. Each query is annotated with source and target types for all three granularity levels; the same query text maps to different type-pairs depending on the ontology. Each predictor receives 10 automatically selected few-shot examples for each granularity level.

\begin{table}[htbp]
  \centering
  \caption{Granularity experiment on K8s subset (17 tools, 128 queries).}
  \label{tab:gran-k8s}
  \small
  \begin{tabular}{llccc}
    \toprule
    & & \textbf{Coarse (6)} & \textbf{Medium (13)} & \textbf{Fine (21)} \\
    \midrule
    \multirow{3}{*}{Oracle}
    & F1            & 0.449 & 0.595          & \textbf{0.706} \\
    & Recall        & 1.000 & 1.000          & 1.000          \\
    & Avg cands.    & 5.1   & 3.7            & \textbf{2.6}   \\
    \midrule
    \multirow{2}{*}{Granite 8B}
    & Exact match   & \textbf{0.773} & 0.750          & 0.711          \\
    & F1        & 0.361 & 0.509          & \textbf{0.586} \\
    \midrule
    \multirow{2}{*}{Opus 4.6}
    & Exact match   & \textbf{1.000} & \textbf{1.000} & \textbf{1.000} \\
    & F1        & 0.449 & 0.595          & \textbf{0.706} \\
    \bottomrule
  \end{tabular}
\end{table}

\begin{table}[htbp]
  \centering
  \caption{Granularity experiment on AAP (1,060 tools, 120 queries). Granite 8B exact match drops to 35\% at 88 types, yet F1 still improves slightly. Opus 4.6 maintains oracle-level performance throughout.}
  \label{tab:gran-aap}
  \small
  \begin{tabular}{llccc}
    \toprule
    & & \textbf{Coarse (8)} & \textbf{Medium (50)} & \textbf{Fine (88)} \\
    \midrule
    \multirow{3}{*}{Oracle}
    & F1            & 0.031  & 0.173          & \textbf{0.230} \\
    & Recall        & 0.958  & \textbf{1.000} & 0.988          \\
    & Avg cands.    & 116.0  & 25.6           & \textbf{21.5}  \\
    \midrule
    \multirow{2}{*}{Granite 8B}
    & Exact match   & 0.608  & \textbf{0.608} & 0.350          \\
    & F1        & 0.022  & 0.110          & \textbf{0.132} \\
    \midrule
    \multirow{2}{*}{Opus 4.6}
    & Exact match   & \textbf{1.000} & \textbf{1.000} & \textbf{1.000} \\
    & F1        & 0.031  & 0.173          & \textbf{0.230} \\
    \bottomrule
  \end{tabular}
\end{table}

\begin{figure}[htbp]
  \centering
  \includegraphics[width=\textwidth]{figures/3_crossover.pdf}
  \caption{K8s (left, 17 tools, 128 queries); AAP (right, 1,060 tools, 120 queries).}
  \label{fig:crossover}
\end{figure}

Tables~\ref{tab:gran-k8s} and~\ref{tab:gran-aap} compare two predictors of different capacity---Granite 4.1 8B and Claude Opus 4.6---against the oracle upper bound. With a perfect predictor (oracle rows), finer ontologies consistently improve the oracle upper bound by reducing candidate ambiguity. The question is whether the predictor can keep up as the type vocabulary grows.

Figure~\ref{fig:crossover} reveals the two opposing effects of ontology granularity. Increasing the number of entity types consistently reduces candidate set size, thereby reducing false positives and increasing routing precision. At the same time, the larger label space makes semantic prediction more difficult for smaller models. On K8s (Figure~\ref{fig:crossover}, left), both predictors improve monotonically with finer types, though the gap between Granite and the oracle widens: exact match drops from 77\% to 71\% as the vocabulary grows from 6 to 21 types. Opus 4.6 matches oracle performance exactly (100\% exact match across all granularities). On AAP (Figure~\ref{fig:crossover}, right), the effect is more pronounced: Granite exact match drops from 61\% to 35\% at 88 types, yet despite substantially lower prediction accuracy, finer ontologies still improve routing precision (F1: 0.110 $\rightarrow$ 0.132) because they generate much smaller candidate sets. Opus 4.6, with sufficient capacity to classify all 88 types correctly, tracks the oracle throughout. This illustrates the trade-off between ontology expressiveness and predictor capacity.

The mechanism is straightforward: finer ontologies reduce the number of tools associated with each source--target type pair ($7.5 \rightarrow 3.6 \rightarrow 3.2$ on AAP), producing smaller candidate sets and fewer false positives. However, each additional entity type must be learned by the predictor. With only 10 few-shot examples per granularity level, Granite struggles to distinguish the 38 fine-grained types added to AAP.

The gap between Granite and Opus 4.6 illustrates the key design trade-off. A sufficiently capable predictor consistently benefits from finer ontologies, whereas weaker predictors eventually become limited by the expanded label space. Ontology granularity is therefore not a fixed property of the API, but a tunable design parameter that should be matched to predictor capacity.

More broadly, ontology granularity controls two opposing effects. Increasing granularity systematically reduces candidate set size, eliminating structurally irrelevant tools and improving routing precision. At the same time, it enlarges the semantic prediction space, making entity-type classification more difficult. The optimal ontology therefore balances these two forces: strong predictors benefit from fine-grained ontologies that maximize precision, while weaker predictors may prefer coarser abstractions that preserve prediction accuracy.

% =============================================================================
% Discussion
% =============================================================================

\section{Discussion}

\textbf{Semantic vs.\ compositional reasoning.} The central finding of this work is that semantic prediction and compositional reasoning have fundamentally different computational properties. TCR exploits this asymmetry by asking the LLM to predict only semantic entity types while delegating compositional reasoning to deterministic graph search. Because the semantic space ($|E|$) is typically much smaller than the tool space ($|T|$), TCR shifts the scaling variable from the number of registered tools to the number of semantic entity types.

An important consequence is that semantic prediction cost is independent of workflow length. Whether a query requires one tool or five, the predictor still outputs only a source and target entity type; additional composition is handled entirely by graph traversal.

The effectiveness of TCR therefore depends critically on the accuracy of entity-type prediction. Deterministic graph search can recover only tool chains consistent with the predicted source and target types; the quality of the candidate set is therefore bounded by the accuracy of semantic prediction. The production-scale AAP experiments illustrate this dependence: Granite 4.1 8B achieves lower routing performance than Opus 4.6 primarily because its lower entity-type prediction accuracy produces less accurate candidate sets. Improving the semantic predictor therefore directly improves end-to-end routing performance without modifying the graph search algorithm.

Likewise, routing quality depends on the completeness of the composition graph. As tool ecosystems evolve, the graph must remain synchronized with tool definitions and type annotations to preserve routing quality.

\textbf{Inference trade-off.} The current implementation replaces one large model invocation over the full tool catalog with two smaller invocations: semantic entity-type prediction followed by tool selection over a reduced candidate set. Although this introduces an additional model invocation, it shifts routing complexity from the number of registered tools to the number of entity types and the size of the candidate set, making the approach increasingly attractive as tool ecosystems scale.

\textbf{The granularity trade-off.} Finer types narrow candidate sets, but the benefit saturates if the predictor cannot keep up with the growing vocabulary. The practical lever is the ratio of candidate-set improvement to classification degradation. With sufficient predictor capacity, finer ontologies consistently improve routing precision by reducing candidate ambiguity.

\textbf{Limitations.}
\begin{itemize}[leftmargin=*,nosep]
  \item \emph{Typed APIs.} TCR assumes tools can be represented as typed transformations between entities.
  \item \emph{Graph maintenance.} Missing or outdated type annotations reduce recall and require the composition graph to remain synchronized with the evolving tool ecosystem.
  \item \emph{Intent ambiguity.} Multiple tools sharing the same type signature still require semantic disambiguation.
  \item \emph{Evaluation.} All benchmark queries were author-annotated; independent benchmarks would strengthen the evaluation. The embedding retrieval baseline (cosine top-10 with nomic-embed-text-v1.5) does not include reranking, BM25, or hybrid search; stronger retrieval baselines may narrow the gap.
\end{itemize}

More broadly, TCR suggests that tool routing should be viewed as the composition of semantic prediction and deterministic graph search rather than as a single end-to-end retrieval problem. Decoupling these two components allows advances in semantic prediction and compositional reasoning to proceed independently: semantic prediction can improve as language models advance, while compositional correctness remains deterministic and verifiable through graph search.

% =============================================================================
% Conclusion
% =============================================================================

\section{Conclusion}

We presented Typed Composition Routing (TCR), a method for tool routing that decouples semantic entity-type prediction from deterministic graph search. Across five benchmark domains and a production-scale API ecosystem with over one thousand tools, this decomposition consistently improved routing quality ($+0.32$ F1 across 20 model--domain combinations), eliminated nonexistent tool invocations by construction, and substantially reduced the prompt footprint required for routing.

We further showed that ontology granularity is itself a key design parameter: finer entity types reduce candidate ambiguity and improve routing precision when supported by sufficiently accurate semantic prediction. This highlights an inherent trade-off between semantic prediction difficulty and candidate ambiguity that should guide the design of future typed tool ecosystems.

By reducing tool routing to classification over a small, fixed vocabulary of entity types (typically tens rather than thousands of classes), TCR also opens the door to lightweight semantic predictors, such as compact encoder models trained specifically for entity-type prediction. Exploring this direction remains future work.

More broadly, this work suggests that improving tool routing may require changing the prediction target: from tools to the domain abstractions that organize them. As tool ecosystems such as Model Context Protocol (MCP) servers and agent platforms continue to grow, we believe typed semantic routing provides a scalable foundation for exposing and selecting increasingly large tool catalogs while dramatically reducing prompt footprint and preserving structural correctness.

% =============================================================================
% Artifact Availability
% =============================================================================

\section*{Artifact Availability}

The implementation of Typed Composition Routing (TCR), benchmark datasets, ontology definitions, evaluation pipeline, experiment configurations, and the production-scale demonstrations described in this paper are available as open-source software. The released artifacts include:

\begin{itemize}[leftmargin=*,nosep]
  \item The TCR routing library;
  \item Benchmark datasets and gold annotations;
  \item Evaluation scripts reproducing all reported metrics and figures;
  \item Ontology definitions used in the experiments;
  \item The production-scale AAP implementation;
  \item The OpenShift MCP server integration demonstrating TCR with real-world tool ecosystems.
\end{itemize}

\noindent Repository URLs:
\begin{itemize}[leftmargin=*,nosep]
  \item \url{https://github.com/jangel97/tcr-go} --- TCR routing library (Go)
  \item \url{https://github.com/jangel97/tcr-k8s-demo} --- K8s benchmark experiments and evaluation pipeline
  \item \url{https://github.com/jangel97/tcr-aap-demo} --- AAP production-scale experiments and evaluation pipeline
  \item \url{https://github.com/jangel97/typed-composition-search} --- Paper, benchmark suite, and evaluation datasets
  \item \url{https://github.com/jangel97/openshift-mcp-server} --- OpenShift MCP server fork with TCR integration
\end{itemize}

\noindent Detailed instructions for reproducing the experiments are provided in the repositories.

% =============================================================================
% References
% =============================================================================

\FloatBarrier
\bibliographystyle{plainnat}
\bibliography{references}

\end{document}
